% !PRINTED 440
\weberpage{440}
%<<CONT
we thus obtain, as the condition that the point $y$ coincide with the point $x$ — since a point is uniquely determined by the \gesperrt{ratios} of its coordinates — the condition that, for a suitably determined coefficient $\lambda$:
\weq{3}{\begin{aligned}
\lambda x_1 &= \alpha\,x_1 + \beta\,x_2 + \gamma\,x_3 \\
\lambda x_2 &= \alpha'x_1 + \beta'x_2 + \gamma'x_3 \\
\lambda x_3 &= \alpha''x_1 + \beta''x_2 + \gamma''x_3,
\end{aligned}}
from which, by elimination of $x_1, x_2, x_3$, one obtains the \gesperrt{cubic} equation for $\lambda$:
\weq{4}{\Lambda = \begin{vmatrix} \alpha - \lambda, & \beta, & \gamma \\ \alpha', & \beta' - \lambda, & \gamma' \\ \alpha'', & \beta'', & \gamma'' - \lambda \end{vmatrix} = 0.}

There are therefore in general \gesperrt{three points} which are their own images, and these we shall call (as with the binary substitutions) the \gesperrt{poles of the substitution $A$}. The three lines joining these poles are then likewise their own images, but in such a way that on each of these three lines there lie only two points which are mapped onto themselves. The images of the remaining points are displaced along the line. These lines we shall call the \gesperrt{axes of the substitution $A$}.

But special cases can now occur which deserve to be noted. Two or even all three poles may coincide in one point, as the example
\[ \begin{pmatrix} \alpha & 0 & 0 \\ \alpha' & \alpha & 0 \\ \alpha'' & \beta'' & \alpha^{-2} \end{pmatrix} \]
shows. Here, if $\alpha'$, $\alpha''$, $\beta''$ are different from zero, two poles coincide in one point, or, if $\alpha = 1$, all three.

More important still is another special case, namely that in which infinitely many points are their own images. If we set aside the case in which all points are their own images, which occurs only in multiplication, then the points which are their own images, if there are infinitely many of them, must lie on a straight line; for this case can occur only if, for a root of (4), the three equations (3) follow from one of them, or, what is the same thing, if together with $\Lambda$ all the first subdeterminants
%CONT>>
